\section{Introduction}\label{intro:section}

Lieb--Thirring inequalities compare the binding energy of a Schr\"odinger
operator with an integral of its potential. Equality asks how the potential
must arrange that binding. In one dimension, a scalar one-bound-state
optimizer has a hyperbolic-secant profile. For a matrix potential, several
such profiles can occupy orthogonal internal channels. The issue is whether
equality also permits channels that rotate with position, or several bound
states that interact within the same channel. We prove that throughout
$1/2<\gamma<3/2$ every equality case consists of independent scalar profiles
in one fixed basis.

Let $m\ge1$ be finite, let $1/2<\gamma<3/2$, and set
$p=\gamma+1/2$. For a measurable Hermitian positive semidefinite function
$W:\R\to\C^{m\times m}$ satisfying
\begin{equation}\label{intro:class}
 \int_\R\tr(W(x)^p)\,dx<\infty,
\end{equation}
let $H_W=-d^2/dx^2-W$ be the operator associated with
\begin{equation}\label{intro:form}
 h_W[\psi]=\int_\R\|\psi'\|^2
       -\int_\R\langle\psi,W\psi\rangle,
 \qquad \psi\in H^1(\R;\C^m).
\end{equation}
Lemma~\ref{pre:spectral} proves that the form is closed and bounded below
and that the negative spectrum is discrete, with possible accumulation
only at zero. We write
\[
 \Tr(H_W)_-^\gamma=\sum_{\lambda_j(H_W)<0}|\lambda_j(H_W)|^\gamma,
\]
counting multiplicities. The finite-dimensional trace $\tr$ is unnormalized;
matrix powers are defined by spectral functional calculus.

\begin{theorem}[Equality classification]\label{thm:main}
For every $W$ satisfying~\eqref{intro:class},
\begin{equation}\label{intro:sharp}
 \Tr(H_W)_-^\gamma\le C_\gamma\int_\R\tr(W^p),
 \qquad
 C_\gamma=\left(\frac{\gamma-1/2}{\gamma+1/2}\right)^{\gamma-1/2}
       \frac{\Gamma(\gamma+1)}{\sqrt\pi\,\Gamma(\gamma+3/2)}.
\end{equation}
Put $r=(\gamma-1/2)^{-1}$. Equality holds if and only if there are
$k\in\{0,\ldots,m\}$, a constant unitary matrix $U$, positive numbers
$a_1,\ldots,a_k$, and real numbers $x_1,\ldots,x_k$ such that almost
everywhere
\begin{equation}\label{intro:classification}
 W(x)=U\diag\bigl(w_1(x),\ldots,w_k(x),0,\ldots,0\bigr)U^*,
 \quad
 w_j(x)=(r+1)a_j^2\sech^2\!\bigl(ra_j(x-x_j)\bigr).
\end{equation}
The case $k=0$ means $W=0$. Each nonzero channel has exactly one negative
eigenvalue, $-a_j^2$, and therefore every extremal potential has at most
$m$ negative eigenvalues, counted with multiplicity.
\end{theorem}

The inequality in~\eqref{intro:sharp} is the sharp matrix inequality
of~\cite[Theorem~1.1]{MLT}. Our contribution is its equality classification.
We retain the analytic constructions in the proof, so that the inequality
and its equality conditions are established here under the same measurable
potential hypothesis. The scales and centers in~\eqref{intro:classification}
are independent; repeated scales and identical profiles are allowed.

\subsection{History and related results}
Lieb and Thirring introduced spectral moment inequalities in their work
on the kinetic energy of fermions and stability of matter
\cite{LiebThirring1975,LiebThirring1976}. Their sharp constants distinguish
the regimes in which binding is best concentrated in a single state or
distributed among many states. Keller's optimization of the lowest
Schr\"odinger eigenvalue~\cite{Keller1961} supplies the one-state problem
underlying the scalar profiles in~\eqref{intro:classification}; their
explicit form appears in~\cite[Section~4(B), Equations~(4.19)--(4.20)]{LiebThirring1976}.
The present question concerns equality in the sum of all negative
eigenvalue moments, with a matrix potential that need not commute with
itself at different points.

The sharp scalar endpoint $\gamma=1/2$ was proved by
Hundertmark, Lieb and Thomas~\cite{HLT1998}; Hundertmark, Laptev and
Weidl~\cite[Theorem~3.1]{HLW2000} extended it to operator-valued
potentials. The semiclassical constant
is sharp for $\gamma\ge3/2$; the lifting of moments of Aizenman and
Lieb~\cite{AizenmanLieb1978} and the operator-valued inequalities of
Laptev and Weidl~\cite{LaptevWeidl2000} are central to this development.
Benguria and Loss~\cite{BenguriaLoss2000} gave a commutation proof of the
matrix result at $\gamma=3/2$. More recently, Read and
Schulz~\cite{ReadSchulz2026} proved the sharp scalar inequality at
$\gamma=1$, with constant $4/(3\sqrt3\pi)$, and its operator-valued
extension. The matrix companion~\cite{MLT} proves the sharp one-state
constant for the whole open interval considered here.

The endpoint $\gamma=3/2$ has a different equality structure. In the
scalar finite-rank problem, Frank, Gontier and Lewin
\cite[Theorem~5]{FrankGontierLewin2025} identify the optimizers as KdV
multisolitons. In particular, several bound states can occur in a single
scalar channel at that endpoint. The conclusion of
Theorem~\ref{thm:main}, one bound state per nonzero channel, uses the
strict inequality $r>1$ and is asserted only for the open interval.

\subsection{Quantitative defects and fixed channels}
The proof begins with the finite-interval matrix action method
of~\cite{MLT}. For orthonormal real trial functions approximating finitely
many negative eigenfunctions, place the functions in the rows of $U(x)$
and choose a lower triangular matrix $C$. The rows of $A=CU$ remain in the
successive trial spans.
If $K$ is the diagonal matrix of positive square roots of the trial
binding energies, the action involves $A'$, $K^2AA^T$, and a power of
the density $AA^T$. A symmetric matrix path $B(x)$ connects $K$ to
$-K$ and gives a lower bound for that action.

For equality, the lower bound must retain more information. Besides the
square $\|A'-BA\|_{\HS}^2$, we retain a term controlling
\[
 [B,K^2]\quad\hbox{and}\quad K^2-B^2-(AA^T)^r.
\]
Its coefficient depends on the largest binding energy and the exponent,
but not on the number of trial functions. This uniformity is the first
essential point: an extremal potential is not known in advance to have
finitely many negative eigenvalues. The integrated trace identity of the
matrix companion supplies the two nonnegative squares from which this
estimate follows.

The second point is compactness with a fixed number of columns. As the
trial list grows, $A$ has more rows but still only $m$ columns. Its
density therefore has rank at most $m$. The retained defect bounds the
tail rows by the small binding energies in that tail. This yields strong
pointwise convergence even for an infinite eigenvalue list and gives
exact algebraic relations in the limit. Those relations separate the
rows with different energies. A differential inequality then makes
their row spaces independent of position; their mutual orthogonality
leaves at most $m$ nonzero directions. On each resulting finite block,
a matrix Riccati equation has a fixed diagonalizing basis and determines
the profiles in~\eqref{intro:classification}.

Section~\ref{pre:section} gives the spectral and convexity preliminaries.
Section~\ref{act:section} proves the quantitative action estimate, and
Section~\ref{bound:section} recovers the sharp spectral bound.
Section~\ref{comp:section} passes to exact relations for an equality
case. Section~\ref{rig:section} proves that the channels are fixed,
solves the Riccati equation, and treats complex potentials and the
converse. Appendices~\ref{app:trace} and~\ref{app:paths} prove the trace
identity, regularized field construction, and connecting-path theorem
used by the action estimate.

\subsection{Notation}
Throughout the proof,
\begin{equation}\label{intro:parameters}
 \sigma=\gamma-\tfrac12\in(0,1),\qquad
 p=1+\sigma,\qquad r=\sigma^{-1},\qquad q=r+1,
\end{equation}
and
\begin{equation}\label{intro:constants}
 D_\sigma=\frac{\sigma^\sigma}{(1+\sigma)^{1+\sigma}},\qquad
 b_\sigma=\int_{-1}^1(1-s^2)^\sigma\,ds,\qquad
 C_\gamma=\frac{D_\sigma}{b_\sigma}.
\end{equation}
The last identity follows from the beta integral. We use the operator
norm $\|\cdot\|_{\op}$ and the Hilbert--Schmidt norm
$\|A\|_{\HS}^2=\tr(A^*A)$. The symbol $T$ denotes transpose of a real
matrix, and $\Sym_n$ is the space of real symmetric $n\times n$ matrices
with inner product $\tr(BH)$. The commutator is $[B,H]=BH-HB$.
The identity matrix, or identity operator, has the size determined by
its context.
